\section{Robustness and Institutional Boundary Calculations}
\label{app:robustness}

This appendix collects the exact calculations underlying the scope statements in Section~\ref{sec:secondary}. They are boundary and falsification results, not additional headline extensions.

\subsection{Zero-network boundary}

Setting $v=0$ in the canonical member welfare gap gives
\begin{equation}
W_M^{SU,N}-W^{IS}
=
\frac{c(13c-6)}{32}.
\label{eq:appendix-zero-network-gap}
\end{equation}
For $0<c<1/3$, $13c-6<13/3-6=-5/3<0$, so the gap is strictly negative.

The zero-network adoption thresholds are
\begin{equation}
T_A(0)=\frac{3c(2-3c)}{16},
\qquad
T_U(0)=T_W(0)=\frac{3c(2-c)}{16}.
\label{eq:appendix-zero-network-thresholds}
\end{equation}
Moreover
\[
2T_A(0)-T_W(0)
=
\frac{3c(2-5c)}{16}>0
\]
for the relevant canonical zero-network range used in the R4 witness analysis, and exact R4 verification identifies positive one-way-adoption intervals even though the initial SU advantage is absent. The full fixed-$c$ root characterization is certified by the Descartes/Sturm calculation in 	exttt{code/revision/check\_r4\_network\_role.py}.

\subsection{Alternative singleton network rule}

R4 pre-specifies one alternative network rule,
\[
g_i=v\sum_{j\in G_i}q_j
\]
for every nonempty compatibility group, including a singleton. Let
\[
K(v)=3v^2-6v+2.
\]
In an SU member market the resulting quantities are
\begin{align}
q_M^{S1}
&=
\frac{1+c-2v}{2K(v)},\
q_B^{S1}
&=
\frac{1-3v-3c(1-v)}{2K(v)},
\end{align}
and in an SU outsider market they are
\begin{align}
q_C^{S1}
&=
\frac{1-2v-2c(1-v)}{2K(v)},\
q_D^{S1}
&=
\frac{1+2c-3v}{2K(v)}.
\end{align}
The old canonical domain keeps all of these quantities strictly positive. The transformed-domain exact sign certificate in the R4 verification script proves
\[
\Phi_{S1}<0
\]
throughout that domain. At the exact point $(c,v)=(1/10,6/25)$, however, the adoption thresholds still satisfy
\[
\max\{T_U^{S1},T_A^{S1},T_W^{S1}\}<2T_A^{S1},
\]
so a positive outsider-only adoption interval remains. The conclusion is therefore specification dependence of the initial political ranking, not failure of private adaptation.

\subsection{Frozen differentiated-demand comparison}

R5 fixes $\gamma=1/2$ and defines, for each compatibility partition,
\[
p=\mathbf 1-M_G(v)q,
\qquad
M_G(v)=B-vC_G,
\]
where
\[
B=(1-\gamma)I+\gamma\mathbf 1\mathbf 1^\top.
\]
Consumer surplus is
\[
CS(q)=\frac12 q^\top M_G(v)q.
\]
The same demand system is solved once with quantity competition and once with price competition, where direct demand is
\[
q=M_G(v)^{-1}(\mathbf 1-p).
\]

For both strategic-variable versions $X\in\{C,B\}$, the exact R5 sign certificates establish
\[
\Phi^X<0,
\qquad
\mathscr E^X<0
\]
throughout the frozen audit box $0<c<1/3$, $0\le v<1/4$ subject to the relevant interiority conditions. At
\[
(c,v,F)
=
\left(
\frac1{10},
\frac6{25},
\frac15
\right),
\]
both versions also satisfy the strict inequalities required for outsider-only adoption. This is the basis for the main-text statement that adoption portability survives while the political sign does not.

\subsection{Strict versus weak blocking}

Let
\[
\operatorname{SB}(x_1,x_2;y_1,y_2)
\]
denote strict blocking by two deviators, requiring $y_k>x_k$ for both $k$, and let weak/Pareto blocking require $y_k\ge x_k$ for both with at least one strict inequality.

Under exact symmetry and high $F$, moving from $SU_{12}$ to $SU_{13}$ yields
\[
W_1^{13,N}=W_1^{12,N},
\qquad
W_3^{13,N}>W_3^{12,N}.
\]
Hence the deviation is not a strict block but is a weak/Pareto block. The same argument applies by permutation to every symmetric SU. The logical distinction is independently formalized in the Lean theorem 	exttt{FV5\_indifferenceSeparatesBlockingRules}.

With
\[
m_1=m_2=1,
\qquad
m_3=1-\delta,
\]
the common-member equality becomes
\[
W_1^{12,N}-W_1^{13,N}
=
\delta(A-C)>0.
\]
The R6 exhaustive deviation enumeration then proves the exact witness ranges stated in Section~\ref{sec:coalition-stability}. In the intermediate full-bypass region, the member and outsider gaps relative to IS are
\[
W_i^{IS}-W_i^{SU,O}
=
m_o(P-C)>0
\]
and
\[
W_o^{IS}-W_o^{SU,O}
=
m_oJ+F>0,
\]
which are strict and therefore do not rely on symmetry-induced indifference.
